\section{Purpose and verification contract}

This blueprint connects the paper
\emph{A Sharper Hoeffding Bound for Weighted Sums of Exchangeable Random
Variables} to named Lean declarations.  The target estimate is
\[
  \E \exp\left\{\lambda \sum_{i=1}^{n}
  w_i(X_i-\overline X_N)\right\}
  \le
  \exp\left\{
    \frac{\lambda^2}{2}\Gamma_N
    \|P_{\mathbf 1^\perp}\widetilde w\|_2^2
  \right\}.
\]

Seongchan Lee leads the Lean formalization for this joint work with his advisor
and coauthor Ilmun Kim.

The development is loaded by the general import
\[
  \texttt{import Testlean.ExchangeableHoeffding.All}.
\]
All declarations use the namespace \texttt{ExchangeableHoeffding}.

The page uses two verification classes.
\begin{itemize}
  \item A \emph{closed proof} has no project-specific analytic-input argument.
  \item A \emph{conditional reduction} is a rigorously checked implication from
  one of the explicitly named analytic-input structures.
\end{itemize}
This distinction is part of the formal specification.  A conditional reduction
is not presented as an end-to-end proof.

\section{Finite vectors and the centered statistic}

\begin{definition}[Finite-vector operations]
\lean{ExchangeableHoeffding.dot, ExchangeableHoeffding.normSq,
ExchangeableHoeffding.avg, ExchangeableHoeffding.projOnePerp,
ExchangeableHoeffding.zeroPad, ExchangeableHoeffding.centeredWeightedSum}
Vectors are functions \(\operatorname{Fin}N\to\R\).  The central definitions are
\[
  \langle a,x\rangle=\sum_i a_i x_i,
  \qquad
  \|a\|_2^2=\sum_i a_i^2,
  \qquad
  \overline x=\frac1N\sum_i x_i,
\]
and
\[
  (P_{\one^\perp}z)_i
  =z_i-\frac1N\sum_jz_j.
\]
The declaration \texttt{zeroPad} embeds the observed weight vector into the
population index set, while \texttt{centeredWeightedSum} is the dot product of
the observation with the projected weight vector.
\end{definition}

\begin{lemma}[Squared norm is nonnegative]
\lean{ExchangeableHoeffding.normSq_nonneg}
For every finite real vector \(a\),
\[
  0\le \|a\|_2^2.
\]
This is a closed finite-sum proof.
\end{lemma}

\begin{lemma}[Measurability and the coordinate-cube bound]
\lean{ExchangeableHoeffding.measurable_dot,
ExchangeableHoeffding.abs_dot_le_sum_abs}
For fixed \(a\), the map \(x\mapsto\langle a,x\rangle\) is measurable.  If
\(|x_i|\le1\) for every coordinate, then
\[
  |\langle a,x\rangle|\le\sum_i|a_i|.
\]
Both are closed proofs and supply the analytic interface used later for
exponential integrability.
\end{lemma}

\begin{lemma}[Zero-padding identities]
\label{lem:zero-pad}
\lean{ExchangeableHoeffding.sum_zeroPad,
ExchangeableHoeffding.dot_zeroPad}
If \(n\le N\), zero-padding preserves both the coordinate sum and the dot
product with the retained coordinates:
\[
  \sum_{i<N}\widetilde w_i=\sum_{i<n}w_i,
  \qquad
  \langle\widetilde w,x\rangle
  =\sum_{i<n}w_i x_i.
\]
Both statements are closed proofs over finite ranges.
\end{lemma}

\begin{theorem}[Projected and coordinate forms agree]
\label{thm:centered-identity}
\lean{ExchangeableHoeffding.centeredWeightedSum_eq_centered_coordinates}
\uses{lem:zero-pad}
The formal statistic agrees exactly with the expression in the paper:
\[
  \texttt{centeredWeightedSum}(w,x)
  =\sum_{i<n} w_i(x_i-\overline x).
\]
The Lean proof expands both finite sums, applies the zero-padding identities,
and closes the remaining real-algebra identity.
\end{theorem}

\begin{lemma}[Projection identities]
\label{lem:projection}
\lean{ExchangeableHoeffding.sum_projOnePerp_eq_zero,
ExchangeableHoeffding.projOnePerp_eq_self,
ExchangeableHoeffding.projOnePerp_idempotent}
For \(N>0\), projection produces a zero-sum vector and is idempotent:
\[
  \sum_i(P_{\one^\perp}z)_i=0,
  \qquad
  P_{\one^\perp}(P_{\one^\perp}z)=P_{\one^\perp}z.
\]
These are closed proofs.
\end{lemma}

\section{Probability model}

\begin{definition}[Exchangeability and boundedness]
\lean{ExchangeableHoeffding.permute,
ExchangeableHoeffding.Exchangeable,
ExchangeableHoeffding.BoundedByOne}
The law \(\mu\) is exchangeable when every coordinate permutation preserves
it under pushforward.  The predicate \texttt{BoundedByOne} records that every
coordinate lies in \([-1,1]\) almost surely.
\end{definition}

\begin{lemma}[Permutation invariance]
\lean{ExchangeableHoeffding.exchangeability_supports_symmetrization}
An exchangeable law is invariant under each named coordinate permutation.
This is the direct formal bridge to the symmetrization step.
\end{lemma}

\section{Inflation factor}

\begin{definition}[Finite-population factor]
\label{def:gamma}
\lean{ExchangeableHoeffding.gammaTerm,
ExchangeableHoeffding.Gamma,
ExchangeableHoeffding.GammaClosed,
ExchangeableHoeffding.harmonic,
ExchangeableHoeffding.epsilonBarber}
The inflation factor is the finite maximum
\[
  \Gamma_N=
  \max_{1\le s\le\lfloor N/2\rfloor}
  \frac{N(N-s)}{s}
  \sum_{\ell=N-s}^{N-1}\frac1{\ell^2}.
\]
The declarations also encode its proposed closed form, the harmonic number, and
the comparison factor \((H_N-1)/(N-H_N)\).
\end{definition}

\begin{lemma}[Endpoint calculation]
\lean{ExchangeableHoeffding.gamma_two,
ExchangeableHoeffding.harmonic_two,
ExchangeableHoeffding.epsilonBarber_two,
ExchangeableHoeffding.gamma_eq_barber_two}
Lean directly checks
\[
  \Gamma_2=2,\qquad H_2=\frac32,\qquad
  \Gamma_2=1+\varepsilon_2.
\]
\end{lemma}

\begin{lemma}[Positivity of the inflation factor]
\lean{ExchangeableHoeffding.gammaTerm_nonneg,
ExchangeableHoeffding.Gamma_nonneg,
ExchangeableHoeffding.gammaTerm_one_pos,
ExchangeableHoeffding.Gamma_pos}
Every candidate is nonnegative, and \(\Gamma_N>0\) for \(N\ge2\).  These are
closed finite-sum proofs.
\end{lemma}

\begin{definition}[General inflation analysis input]
\lean{ExchangeableHoeffding.GammaAnalyticInputs}
The closed formula, asymptotic expansion, and strict comparison for all larger
populations are grouped in \texttt{GammaAnalyticInputs}.  The declaration
\texttt{gamma\_strictly\_improves\_barber} is therefore a conditional
reduction from this input, not yet a closed theorem.
\end{definition}

\section{Slice reduction}

\begin{definition}[Hamming-slice quantities]
\lean{ExchangeableHoeffding.sliceSubsets,
ExchangeableHoeffding.sliceAverage,
ExchangeableHoeffding.elemSym,
ExchangeableHoeffding.sliceFunctional,
ExchangeableHoeffding.AtMostTwoValues,
ExchangeableHoeffding.SphereSection}
For a uniformly selected \(k\)-subset \(S\), the target slice estimate is
\[
  \log\E_S\exp\left\{\sum_{i\in S}y_i\right\}
  \le \frac{\Gamma_N}{8}\|y\|_2^2,
  \qquad \sum_i y_i=0.
\]
The functional \texttt{sliceFunctional} is the corresponding logarithm of an
elementary symmetric polynomial normalized by \(\binom Nk\).
\end{definition}

\begin{lemma}[Slice expectation identity]
\lean{ExchangeableHoeffding.sliceAverage_exp_sum_eq_elemSym,
ExchangeableHoeffding.sliceFunctional_eq_log_sliceAverage}
Lean proves that the exponential average over the finite slice is exactly the
normalized elementary symmetric polynomial.  Thus the probabilistic and
variational definitions of the slice log-MGF coincide.
\end{lemma}

\begin{lemma}[Finite slice and two-level algebra]
\lean{ExchangeableHoeffding.sliceSubsets_eq_powersetCard,
ExchangeableHoeffding.card_sliceSubsets,
ExchangeableHoeffding.sliceAverage_one,
ExchangeableHoeffding.centeredTwoLevelVector,
ExchangeableHoeffding.sum_centeredTwoLevelVector,
ExchangeableHoeffding.normSq_centeredTwoLevelVector,
ExchangeableHoeffding.sum_centeredTwoLevelVector_on}
The \(k\)-slice has cardinality \(\binom Nk\), and its normalized average sends
the constant one function to one.  Lean also checks the paper's centered
two-level parameterization, its squared norm, and the identity reducing a slice
sum to a centered intersection count.
\end{lemma}

\begin{definition}[Three-coordinate extremizer input]
\lean{ExchangeableHoeffding.threeObjective,
ExchangeableHoeffding.ThreeSphereSection,
ExchangeableHoeffding.HasDuplicateCoordinate,
ExchangeableHoeffding.ThreePointAnalyticInputs}
The three-coordinate stage studies
\[
  A\sum_{i=1}^3e^{x_i}+B\sum_{i=1}^3e^{-x_i}
\]
under fixed coordinate sum and squared norm.  The derivative-sign statement and
the repeated-coordinate conclusion are presently explicit analytic inputs.
The two reciprocal-Lagrange-weight tangent identities used by the second-order
argument are already closed as
\texttt{lagrange\_tangent\_sum\_zero} and
\texttt{lagrange\_tangent\_weighted\_sum\_zero}.
\end{definition}

\begin{definition}[Two-level slice input]
\lean{ExchangeableHoeffding.SliceAnalyticInputs}
The two-level extremizer and the final Hamming-slice MGF estimate are fields of
\texttt{SliceAnalyticInputs}.  Closing this structure requires turning the
three-coordinate reduction into a global extremizer proof and then evaluating
the resulting two-level vector.
\end{definition}

\section{Hypergeometric reduction}

\begin{definition}[Sampling quantities]
\lean{ExchangeableHoeffding.hypergeomPMF,
ExchangeableHoeffding.hypergeomMgf,
ExchangeableHoeffding.hypergeomB}
After the two-level reduction, the slice sum becomes a centered
hypergeometric count.  Its finite MGF is \texttt{hypergeomMgf}, and the target
quadratic coefficient is
\[
  B_{N,m}=(N-s)^2\sum_{\ell=N-s}^{N-1}\ell^{-2},
  \qquad s=\min(m,N-m).
\]
\end{definition}

\begin{definition}[Hypergeometric analytic input]
\lean{ExchangeableHoeffding.HypergeometricAnalyticInputs}
This structure contains only the project-specific sharpened hypergeometric estimate
\[
  \log\E e^{t(H-mK/N)}\le \frac{t^2}{8}B_{N,m}.
\]
It remains an explicit completion target.
\end{definition}

\begin{lemma}[Bounded centered MGF]
\lean{ExchangeableHoeffding.bounded_centered_mgf}
The ordinary Hoeffding lemma in the normalization needed by the martingale
argument is a closed theorem.  Its formal statement makes measurability,
centering, and almost-sure interval bounds explicit.
\end{lemma}

\section{Main results and dependency map}

\begin{definition}[Main analytic input]
\label{def:main-input}
\lean{ExchangeableHoeffding.MainAnalyticInputs}
The main input packages two paper-level statements: the zero-sum MGF bound and
the sharp lower bound for admissible constants.  The one-sided tail estimate is
derived from the MGF statement and is not an input field.
\end{definition}

\begin{theorem}[Centered exchangeable MGF reduction]
\label{thm:mgf}
\lean{ExchangeableHoeffding.exchangeable_mgf_bound}
\uses{lem:projection, def:gamma, def:main-input}
Given \texttt{MainAnalyticInputs}, Lean proves the centered MGF estimate by
setting \(a=P_{\one^\perp}\widetilde w\), invoking the closed zero-sum
projection identity, and applying the zero-sum MGF field.  The reduction itself
is checked; the analytic input is not constructed in the repository.
\end{theorem}

\begin{corollary}[One-sided tail reduction]
\lean{ExchangeableHoeffding.tailThreshold,
ExchangeableHoeffding.integrable_exp_mul_dot_of_bounded,
ExchangeableHoeffding.exchangeable_tail_bound_from_mgf,
ExchangeableHoeffding.exchangeable_tail_bound}
The exported threshold is
\[
  \|P_{\one^\perp}\widetilde w\|_2
  \sqrt{2\Gamma_N\log(1/\delta)}.
\]
Lean proves the Chernoff step from the MGF field: bounded coordinates give
exponential integrability, positivity of \(\Gamma_N\) makes the sub-Gaussian
parameter valid, and the threshold algebra reduces the exponential bound to
\(\delta\).  The result remains conditional only through the zero-sum MGF field.
\end{corollary}

\begin{proposition}[Sharp admissible-constant reduction]
\lean{ExchangeableHoeffding.AdmissibleConstant,
ExchangeableHoeffding.optimalLowerBound,
ExchangeableHoeffding.admissible_constant_lower_bound}
Any uniformly admissible replacement constant is intended to satisfy
\[
  C\ge
  \begin{cases}
  N/(N-1),&N\text{ even},\\
  (N+1)/N,&N\text{ odd}.
  \end{cases}
\]
The current Lean declaration is a checked export of the corresponding main
analytic input.
\end{proposition}

\section{Completion criterion}

The finite-algebra bridge through Theorem~\ref{thm:centered-identity} is closed.
The end-to-end probability theorem is complete only after Lean constructors are
proved for
\[
\begin{gathered}
  \texttt{GammaAnalyticInputs},\quad
  \texttt{HypergeometricAnalyticInputs},\quad
  \texttt{ThreePointAnalyticInputs},\\
  \texttt{SliceAnalyticInputs},\quad
  \texttt{MainAnalyticInputs}.
\end{gathered}
\]
At that point the public MGF, tail, and sharpness theorems can drop their input
arguments.  Until then, the project page records the MGF dependency explicitly
while marking the Hoeffding and Chernoff steps as closed.
